\documentclass[10pt,a4paper]{amsart}
\usepackage[margin=19mm]{geometry}
\usepackage{amsmath,amssymb,mathtools}
\usepackage[scr=rsfs,cal=euler]{mathalfa}
\usepackage{xfrac}
\usepackage[unicode,hidelinks,bookmarks=false]{hyperref}
\newcommand{\ie}{i.e.\ }
\newcommand{\cf}{cf.\ }
\newcommand{\resp}{respectively\ }
\newcommand{\Subsection}{\subsection}
\providecommand{\nobreakdash}{\nobreak\hbox{-}}
\setlength{\emergencystretch}{2em}
\multlinegap0pt
\usepackage{amssymb}
\usepackage{mathtools}
\usepackage{float}
\newtheorem{theorem}{Theorem}
\newcommand{\bA}{\mathbf{A}}
\newcommand{\bW}{\mathbf{W}}
\newcommand{\br}{\mathbf{r}}
\newcommand{\valpha}{{\boldsymbol{\alpha}}}
\newcommand{\vphi}{\boldsymbol{\phi}}
\newcommand{\vpsi}{\boldsymbol{\psi}}
\title{Generalized Transferable Neural Networks for Steady-State Partial Differential Equations}
\author{Tao Cheng, Lili Ju, Zhonghua Qiao, Xiaoping Zhang}
\date{Source: arXiv:2604.03020v1 (CC BY 4.0)}
\begin{document}
\maketitle

\noindent\emph{Source and licence.} Generalized Transferable Neural Networks for Steady-State Partial Differential Equations. Version arXiv:2604.03020v1 (CC BY 4.0), distributed by arXiv under the Creative Commons Attribution 4.0 International licence, http://creativecommons.org/licenses/by/4.0/, as recorded in the arXiv licence metadata for this exact version and shown on the abstract page https://arxiv.org/abs/2604.03020v1. Full licence notice: \texttt{licenses/arxiv-2604.03020-CC-BY-4.0.txt}.\par\smallskip

\noindent\textit{Editorial scope.} Counted unit: the article's theorem thm:1 (source0.tex lines 551--558). The article's complete original proof of it is delivered with it (source0.tex lines 560--611, one \texttt{proof} environment of 52 lines). No mathematics is written, repaired or re-derived: the statement and the proof are the article's own lines, in the article's own notation, and no retained line is substituted or re-flowed.  Retained uncounted as the source-local context the proof needs: lines 431--439; lines 506--510; lines 520--522; lines 524--526.  Nothing was omitted from any retained span, so the package carries no omission record.  No retained line contains a citation, so the wrapper carries no bibliography and no bibliography entry.  No retained line contains a figure environment, an image-inclusion command or drawing code, so no image is delivered and the package declares no asset dependency.  Editorial formatting only, in addition: an AMS \texttt{amsart} preamble that restates the article's own declarations and macros used by the retained lines, namely the environments \texttt{theorem} on separate counters, together with the standard packages those lines need.  The retained environments print in this document as Theorem 1 in order of appearance, whereas the article prints the same items with the numbering of its own full document.  The complete native source of the article as distributed in the arXiv e-print for this exact version is bundled unchanged as \texttt{sources/197780/source0.tex}.
\begin{equation}\label{eq:g_transnet}
    \begin{cases}
        \vpsi_1 = \tanh({\mathbf\Gamma}\odot(\bA(\mathbf{x}-\mathbf{x}_c) + R\br)), \\
        \vpsi_2 = \tanh(\bW_2\vpsi_1), \\
        \vdots\\
        \vpsi_L = \tanh(\bW_L\vpsi_{L-1}), \\
        u_{\mathrm{NN}}(\mathbf{x}) = \valpha^\top \vpsi_L,
    \end{cases}
\end{equation}

\begin{equation}\label{eq:sample_W}
W_{ij}^{(l)} \sim \mathcal{N} \left(0,\sigma_l^2\right)
\quad{\rm with}\;\;
\sigma_l = \textstyle\sqrt{\frac{\delta}{N_{l-1}}},
\end{equation}

\begin{equation}\label{eq:E_phi}
\mathbb{E}[\vphi_l] = \mathbf{0}, \qquad l=2,\dots,L,
\end{equation}

\begin{equation}\label{eq:E_psi}
\mathbb{E}[\vpsi_l] = \mathbf{0}, \qquad l=2,\dots,L.
\end{equation}

\begin{theorem}[Controlled Variance Propagation in GTransNet]\label{thm:1}
Assume that each component of $\vpsi_1$ has a variance not more than $\sigma_0^2$. 
Then the variances of the neuron activations  in  the subsequent hidden layers of GTransNet \eqref{eq:g_transnet}  satisfy
\begin{equation}\label{eq:var_psi}
\mathrm{Var}[\psi_i^{(l)}] \le \delta^{\,l-1}\sigma_0^2,
\qquad i=1,\dots,N_l,\quad l=2,\dots,L.
\end{equation} 
\end{theorem}

\begin{proof}


Combining  \eqref{eq:E_phi}-\eqref{eq:E_psi} with the independence of $\bW_l$ and $\vpsi_{l-1}$, and by noting $\mathbb{E}\left[W_{ij}^{(l)}\right] = \mathbb{E}\left[\psi_{j}^{(l-1)}\right] = 0$,  we get 
\begin{equation}\label{eq:var_phi} 
   \begin{aligned}
       \mathrm{Var}\left[\phi_i^{(l)}\right] &= \mathbb{E}\left[\left(\phi_i^{(l)}\right)^2\right] - \mathbb{E}\left[\phi_i^{(l)}\right]^2 
       = \textstyle\mathbb{E}\left[\left(\sum\limits_{j=1}^{N_{l-1}} W_{ij}^{(l)}\psi_{j}^{(l-1)}\right)^2\right] \\
       &=\textstyle \mathbb{E}\left[\sum\limits_{j=1}^{N_{l-1}} \left(W_{ij}^{(l)}\right)^2 \left(\psi_{j}^{(l-1)}\right)^2\right] 
       = \textstyle\sum\limits_{j=1}^{N_{l-1}} \mathbb{E}\left[\left(W_{ij}^{(l)}\right)^2\right] \mathbb{E}\left[\left(\psi_{j}^{(l-1)}\right)^2\right]\\
       &= \sum_{j=1}^{N_{l-1}}\mathrm{Var}\left[W_{ij}^{(l)}\right]\mathrm{Var}\left[\psi_j^{(l-1)}\right],
   \end{aligned}
   \end{equation}
for  $i=1,\dots,N_l$ and $l=2,\dots,L$.
If each component of $\vpsi_1$ has a variance not more than $\sigma_0^2$, then for the second hidden layer we have from \eqref{eq:sample_W} that
\begin{equation}
\mathrm{Var}\left[\phi_i^{(2)}\right] 
= \textstyle\sum\limits_{j=1}^{N_1}\mathrm{Var}\left[W_{ij}^{(2)}\right]\mathrm{Var}\left[\psi_j^{(1)}\right]
= N_{1} \sigma_2^2 \mathrm{Var}\left[\psi_{j}^{(1)}\right]
        \leq \delta \sigma_0^2.
\end{equation}
Using the fact that $|\tanh(x)| \le |x|$ for any $x\in{\mathbf R}$ and $\mathbb{E}[\vpsi_l]=0$, we obtain 
\begin{equation}\label{eq:var_psi1}
\begin{aligned}
\mathrm{Var}\left[\psi_i^{(2)}\right]
&= \mathbb{E}\left[\left(\psi_i^{(2)}\right)^2\right]
= \mathbb{E}\left[\tanh^2\left(\phi_i^{(2)}\right)\right] \\
&\le \mathbb{E}\left[\left(\phi_i^{(2)}\right)^2\right]
= \mathrm{Var}\left[\phi_i^{(2)}\right] 
\leq \delta \sigma_0^2.
\end{aligned}
\end{equation}
Similarly, for the third hidden layer we have from 
\eqref{eq:var_psi1} that
\begin{equation}
\mathrm{Var}\left[\phi_i^{(3)}\right]
= \textstyle\sum\limits_{j=1}^{N_2}\mathrm{Var}\left[W_{ij}^{(3)}\right]\mathrm{Var}\left[\psi_j^{(2)}\right]
= N_2 \sigma_3^2 \mathrm{Var}\left[\psi_i^{(2)}\right]
\le \delta(\delta \sigma_0^2)= \delta^2 \sigma_0^2
\end{equation}
and consequently
\begin{equation}
\begin{aligned}
\mathrm{Var}\left[\psi_i^{(3)}\right]
&= \mathbb{E}\left[\left(\psi_i^{(2)}\right)^2\right]
\le \mathbb{E}\left[\left(\phi_i^{(2)}\right)^2\right]
= \mathrm{Var}\left[\phi_i^{(2)}\right] 
\leq \delta^2 \sigma_0^2.
\end{aligned}
\end{equation}
By repeating this argument sequentially across all  hidden layers in the same manner, we arrive at the final result \eqref{eq:var_psi} and the proof is completed.
\end{proof}

\end{document}
